% HD-S45. GPT-6 Astra Pro /root/sixfolds; inherited effort not exposed.
% 21 November 2026. Exact remaining multiplicity-three tensor source.
\subsection{HD-S45: the multiplicity-three tensor source}
\label{sec:multiplicity-three-tensor}

Work over \(K=\mathbb Q(i)\) or
\(K=\mathbb Q(\sqrt{-d})\), \(d=2,7,11,19,43,67,163\),
with \(K\) the exact coefficient character field. In the
multiplicity-three case of HD-S44, every natural sixfold
source of period rank nine has an already established
determinant class, except possibly for the following exact
remaining form:
\[
 H=S_s\times\Gamma,\qquad
 V=\operatorname{Aug}(K^s)\otimes_K M,\qquad
 s=r+1,\quad r\ge9,\quad \dim_KM=3.                         \tag{S45.1}
\]
Here \(\Gamma\subset GL_K(M)\) is a finite absolutely
irreducible group with exact character field \(K\).
The product, all coefficient maps and its actual Jacobian
quotient description are proved below. No assertion of
dominance or new determinant coverage for (S45.1) is made.
If \(\Gamma\) has a \(K\)-defined line system, the remaining
form also has an already established determinant class.

\paragraph{1. The actual tensor system and its rank bound.}
Use the literal \(AP\otimes_K PU\) map (S44.9), and
then the component quotient maps (S44.13)--(S44.14).
The resulting faithful coefficient group acts on
\[
                   U=W\otimes_K M,\quad
                   \dim W=r,\quad \dim M=3,                 \tag{S45.2}
\]
and contains an actual normal irreducible reflection group
\(L\) on \(W\), acting identically on \(M\).
Its tensor-factor projective maps are exactly those obtained
from transformed matrix units in HD-S44. A prescribed original
primitive line is carried by the displayed matrix-coefficient
isogeny; all polarization multipliers have norm-valued sixth
power. Thus the determinant and period conclusions below
apply to that original factor.

The complete HD-S36, HD-S41--HD-S43 and HD-G30 reductions
already cover every coefficient group with a nonidentity
element of codimension one or two. It remains to analyze
the groups in (S45.2) for which
\[
                    \operatorname{rank}(h-I)\ge3
                         \quad(h\ne I).                    \tag{S45.3}
\]
At period rank nine the quotient base is \(\mathbb P^1\):
genus one has at most six marked parameters, and genus at
least two is excluded by the Euler formula, except for
rank three in genus two where the dimension is three.
Consequently the true inertia tuple has
\[
             b\ge12,\qquad \sum_{j=1}^b c_j=6r+6.           \tag{S45.4}
\]

The projective image on \(M\) is not cyclic. Otherwise
all its matrices would be scalar multiples of powers of
one semisimple matrix, preserving an eigenline; tensoring
that eigenline with \(W\) would contradict irreducibility.
The projective product-one tuple therefore has at least
three nonidentity entries. For any such entry, a nonscalar
matrix on \(M\) has maximum eigenvalue multiplicity at
most two. The fixed dimension on \(W\otimes M\) is
therefore at most \(2r\), so its codimension is at least
\(r\). Every other nonidentity inertia has codimension at
least three by (S45.3). If \(r\ge3\), this gives
\[
                 6r+6\ge3r+3(b-3)\ge3r+27,
                 \qquad r\ge7.                            \tag{S45.5}
\]
The smaller ranks have coefficient rank at most six and
are already covered by HD-S40. Thus no case has been
dropped in deriving \(r\ge7\).

Cohen's classification at ranks at least seven consists
of the real rational types and the imprimitive groups.
The precise list and its matrix-field verification are
the ones in HD-S44, paragraph3. The only nonimprimitive
possibilities are \(A_r\) and, at their indicated ranks,
\(E_7,E_8\).

\paragraph{2. Imprimitive factors and a scalar-product lemma.}
For \(L=G(\ell,e,r)\), the exact trace recovery (S44.10)
gives \(\zeta_\ell\in K\). The characteristic diagonal
subgroup has actual \(K\)-rank-one coordinate projectors
on \(W\). Tensoring their images with \(M\) yields
orthogonal \(K\)-rank-three blocks for \(H\).
Its block stabilizer is absolutely irreducible, because a
proper invariant subspace would induce one for \(H\).
The induced functions, quotient, corner and polarization maps
are literally (S44.13)--(S44.14). They identify the source
factor with an actual faithful coefficient model of rank
three and exact field \(K\). At period rank nine its
base is \(\mathbb P^1\), with twelve reflection inertia.
HD-S40 therefore gives an already established determinant
class. This argument retains the original primitive line.

We will use the following elementary fact, with its exact
linear proof. An absolutely irreducible representation of
dimension \(q>1\) cannot be generated by \(q\) reflections
\(R_1,\ldots,R_q\) whose product is scalar. Indeed the
sum of their rank-one images is invariant under each
generator: if \(u_i\in\operatorname{im}(R_i-I)\), then
\(R_j u_i-u_i\in\operatorname{im}(R_j-I)\).
Irreducibility makes this sum the whole space, so the
\(q\) image vectors form a basis. In that basis \(R_j-I\)
has only row \(j\) nonzero. If
\[
                       R_1\cdots R_q=\lambda I,             \tag{S45.6}
\]
the last row on the left is the last row of \(R_q\).
It follows that \(R_q\) is diagonal, with nontrivial
entry \(\lambda\) in its last coordinate. Removing it,
the last row of \(R_1\cdots R_{q-1}\) is the last
standard row and all other target diagonal entries are
\(\lambda\). Repeating the same row argument makes every
\(R_j\) diagonal, with entry \(\lambda\) in its own
coordinate. If \(\lambda=1\) these were not reflections;
otherwise the representation is reducible. Both alternatives
contradict the assumptions. No finite-group classification
is used in this lemma.

\paragraph{3. The real factors of ranks seven and eight.}
The remaining real groups at these ranks are
\(A_7,A_8,E_7,E_8\). Their full complex linear
normalizers are scalar times their rational Weyl models.
For \(A_r\) this is the explicit root-star or cubic-Hessian
normalizer proved in HD-S36 and HD-S41. For \(E_7,E_8\),
the chamber map of HD-S36 sends an orthogonal root normalizer
to a simple-root diagram automorphism; both diagrams have
trivial automorphism group. The primary locators for this
actual map are Humphreys, \emph{Introduction to Lie Algebras
and Representation Theory}, Theorem10.3(a),(b),(e),
pp.51--52, and Section12.2 with Table1, pp.65--66.
The passage from a complex linear normalizer to a scalar
times a real orthogonal map follows from uniqueness of the
invariant quadratic form and the \(g^{-1}\bar g\) scalar
calculation, as written in HD-S36. All rational Weyl models
are transferred by the finite \(K\)-linear intertwiner.

Adjoin \(-I\) to the first Weyl group, writing \(L^\pm\).
For every inertia element choose the actual tensor lift
\[
                   h_j=A_j\otimes B_j,\qquad
                   A_j\in L^\pm,\quad B_j\in GL_3(K).
                                                               \tag{S45.7}
\]
The scalar in the first normalizer belongs to \(K\):
compare one nonzero entry of its rational Weyl matrix with
the original \(K\)-matrix. It is absorbed into \(B_j\).
The inverse image of the finite group \(H\) under the
tensor map \(L^\pm\times GL_3(K)\to GL(U)\) is finite,
since its kernel is exactly
\(\{(I,I),(-I,-I)\}\). Hence all \(B_j\) lie in a
finite group and are semisimple. Their generated group is
absolutely irreducible; an invariant subspace would tensor
to an \(H\)-invariant subspace. The product-one relation
gives the exact lifted identities
\[
            \prod A_j=\epsilon I,\qquad
            \prod B_j=\epsilon I,\qquad \epsilon\in\{1,-1\}.
                                                               \tag{S45.8}
\]
Thus no unjustified choice of a product-one tensor lift
has been made.

There are exactly three projectively nontrivial \(B_j\).
Four would give \(4r+3(b-4)\ge4r+24>6r+6\) for
\(r=7,8\). The minimal contribution of three is
\(3r+3(b-3)\); the total slack is at most zero for
\(r=7\), and at most three for \(r=8\).

For a projectively trivial \(B_j=aI\), a nonreal \(a\)
would require a nonreal eigenvalue of the real finite-order
matrix \(A_j\) in its fixed tensor space. Such an eigenvalue
has multiplicity at most \(\lfloor r/2\rfloor\), since
its conjugate has the same multiplicity. The codimension
would be at least \(3\lceil r/2\rceil\ge12\), exceeding
its baseline three by at least nine. This is impossible
within the stated slack. Therefore \(a=\pm1\), and
absorbing the same central sign into \(A_j\) makes \(B_j=I\).

For each of the three projectively nontrivial matrices, three
distinct eigenvalues would give fixed dimension at most \(r\),
and codimension at least \(2r\), exceeding the baseline \(r\)
by more than the slack. Its eigenvalues therefore have
multiplicities two and one. Let \(\beta\) be the double
eigenvalue. If \(\beta\) were nonreal, the multiplicity \(u\)
of \(\beta^{-1}\) in the real \(A_j\) would be at most
\(\lfloor r/2\rfloor\). The other matched eigenvalue has
multiplicity at most \(r-u\). The fixed tensor dimension
would be at most \(2u+(r-u)\le r+\lfloor r/2\rfloor\).
Its codimension would exceed \(r\) by at least
\(\lceil r/2\rceil\ge4\), again impossible.
Thus \(\beta=\pm1\). Absorb this central sign into \(A_j\);
the new \(B_j\) fixes a plane and is a reflection.

All actual tensor elements remain unchanged by these sign
absorptions, so their \(B_j\)'s still generate an absolutely
irreducible group. They are now exactly three reflections,
with scalar product (S45.8). The \(q=3\) instance of
(S45.6) contradicts irreducibility. This excludes all four
real rank-seven and rank-eight cases. The argument also
works over any imaginary quadratic field; the determinant
scope restriction entered only in the earlier covered cases.

\paragraph{4. The remaining actual product and its quotients.}
It remains that \(r\ge9\) and \(L=A_r\).
After the actual \(K\)-intertwiner,
\(W=\operatorname{Aug}(K^s)\), \(s=r+1\).
The normalizer just proved in HD-S41 makes every tensor
element uniquely
\[
                         P_\pi|_W\otimes B.
 \tag{S45.9}
\]
Uniqueness follows because the only scalar permutation on
the augmentation space is the identity. These coordinates
are therefore an actual injective homomorphism
\(H\to S_s\times GL_3(K)\). The supported subgroup
contains \(S_s\times1\); multiplying a pair \((\pi,B)\)
by \((\pi^{-1},1)\) proves \(H=S_s\times\Gamma\)
literally. The group \(\Gamma\) is finite and absolutely
irreducible. Since \(\operatorname{tr}(1,B)=r\operatorname{tr}B\)
and every permutation augmentation trace is rational, its
exact character field is \(K\). Schur invariance gives
the actual coefficient form \(h_W\otimes h_M\):
fix \(w\ne0\), define \(h_M\) by restriction to
\(w\otimes M\) divided by \(h_W(w,w)\), and uniqueness
of the \(S_s\)-invariant form verifies the tensor identity.

Here is the remaining source in actual curve and Jacobian
maps. Let \(C\) be its connected product-group cover,
\(L_0=S_{s-1}\), and put
\[
 \begin{gathered}
 X=C/(L_0\times1),\qquad Z=C/(S_s\times1),\qquad
 Y=C/(L_0\times\Gamma),\\
 \alpha:C\longrightarrow X,\qquad f:X\longrightarrow Z,
                         \qquad\deg f=s.                   \tag{S45.10}
 \end{gathered}
\]
The subgroup intersection and product identities identify
\(X\) with the normalization of \(Y\times_{\mathbb P^1}Z\);
both \(X/Y\) and \(Z/\mathbb P^1\) have full group \(\Gamma\).
For its rank-three system \(\mathcal M\), the actual
diagonal inclusion \(u\mapsto(u,\ldots,u)\) and trace
\(s^{-1}\sum u_i\) split
\[
 H^1(Y,j_*\mathcal M_Y)=
 H^1(\mathbb P^1,j_*\mathcal M)\ \perp\
 H^1(\mathbb P^1,j_*(W\otimes\mathcal M)).
                                                               \tag{S45.11}
\]
The first summand's coefficient form is multiplied by \(s\).
The local stalk map on a permutation cycle of length \(\ell\)
uses the product monodromy \(\gamma^\ell\); hence this
splitting includes all ramified stalks.

Write \(d_0=\dim_KH^1(\mathbb P^1,j_*\mathcal M)\), and
let \(e_u\in\mathbb Q[\Gamma]\) be the orthogonal primitive
corner for a chosen \(K\)-line, including its conjugate.
The sixfold is the connected image on \(JX\) of an integer
multiple of the actual self-adjoint projector
\[
 \pi=e_u\left(1-\frac1s f^*f_*\right),\qquad
 \frac1{|L_0|}\alpha^*\pi\alpha_*
            =(\operatorname{av}_{L_0}-\operatorname{av}_{S_s})e_u.
                                                               \tag{S45.12}
\]
These identities follow from \(f_*f^*=[s]\) and the two
quotient norm identities, exactly as in HD-S43.
Thus it is the polarized complement of the actual pullback
of the \(d_0\)-dimensional factor of \(JZ\) inside the
\((6+d_0)\)-dimensional \(e_u\)-factor of \(JX\).
This displayed factor of \(JZ\) is allowed to vary here;
no constancy or period bound for this remaining construction
has been asserted.

Finally suppose \(\Gamma\) has a \(K\)-defined line system.
For a line stabilizer \(J_\Gamma\), its character \(\chi\)
has values in the roots of unity of \(K\). The original
coefficient representation is then literally
\[
 \operatorname{Ind}_{S_s\times J_\Gamma}^{S_s\times\Gamma}
                      (W\otimes\chi),                     \tag{S45.13}
\]
with the same coefficient-vector, quotient and polarization
maps (S44.13)--(S44.14). Its faithful block image is
\(S_s\times\chi(J_\Gamma)\). Its coefficient rank is
\(r\ge9\), so a rank-nine family again has genus-zero
quotient base: genus one has at most six parameters and
the Euler formula excludes higher genus. In the Gaussian case,
exact field \(K\) forces \(\chi(J_\Gamma)=\mu_4\);
HD-S33 applies to this actual scalar-augmentation source.
For the seven other fields the character has values
\(\{\pm1\}\), so the induced trace would be rational,
contradicting the exact field \(K\). This proves the
last stated reduction while retaining the specified corner.
